\documentclass[11pt]{article}
\usepackage[margin=1in]{geometry}
\usepackage{amsmath,amssymb,amsthm,mathtools}
\usepackage{algorithm,algpseudocode}
\usepackage{natbib}
\usepackage{booktabs}
\usepackage{hyperref}
\usepackage{cleveref}
\usepackage{bm}
\usepackage{microtype}
\usepackage{pgfplots,pgfplotstable}
\pgfplotsset{compat=1.18}
\usepgfplotslibrary{groupplots}
\usepackage{xcolor}
\usepackage{subcaption}
\usepackage{multirow}

\definecolor{colcases}{HTML}{CC0000}
\definecolor{colhosp}{HTML}{7B2D8E}
\definecolor{colicu}{HTML}{8B0000}
\definecolor{colradar}{HTML}{2E8B57}
\definecolor{colmodel}{HTML}{006400}
\definecolor{colrivm}{HTML}{666666}
\definecolor{colcori}{HTML}{FF8C00}

\newcommand{\bx}{\bm{x}}
\newcommand{\bz}{\bm{z}}
\newcommand{\bc}{\bm{c}}
\newcommand{\ba}{\bm{a}}
\newcommand{\btheta}{\bm{\theta}}
\DeclareMathOperator*{\argmin}{arg\,min}
\newcommand{\dd}{\,\mathrm{d}}
\newcommand{\doi}[1]{\href{https://doi.org/#1}{\texttt{doi:#1}}}
\newtheorem{remark}{Remark}

\title{Calibration of Epidemic Models via Multi-Source Generalized Profiling}
\author{Hadi A.\ Peivasti \and Martijn Schoot Uiterkamp\\[4pt]
\small Department of Econometrics and Operations Research\\
\small Tilburg University, The Netherlands}
\date{March 2026}

\begin{document}
\maketitle

\begin{abstract}
Calibrating compartmental epidemic models to surveillance data is difficult when multiple noisy and indirect data streams are available and the transmission rate varies over time. We develop a generalized profiling framework for SEIR models that fuses heterogeneous data sources---reported cases, hospital and ICU admissions, and digital contact monitoring data from the COVID~RADAR app---through epidemiologically motivated observation operators. State trajectories are represented as log-transformed B-splines, which guarantee positivity without inequality constraints, and are estimated jointly with a spline-based time-varying transmission rate via parameter cascading. Applied to Dutch COVID-19 surveillance data from the first pandemic year (March 2020--March 2021), the method recovers a transmission rate whose implied reproduction number correlates at $r=0.82$ (Spearman $r_s=0.83$, $R^2=0.54$) with independently estimated values from the Dutch National Institute for Public Health and the Environment (RIVM), and agrees with RIVM on the direction of epidemic growth ($R_t>1$ vs.\ $R_t<1$) on 75\% of days, despite these estimates not being used in fitting. A fully independent prevalence validation achieves $r=0.85$ with zero shared input data. A source ablation study across eight data configurations reveals that digital contact monitoring data provides the highest marginal information value among all streams: fitting hospital admissions, ICU data, and RADAR contact data---without any case data---achieves the highest prevalence correlation ($0.90$) of all configurations tested. A systematic weight optimization over 64~configurations and a penalty weight sensitivity analysis further demonstrate that the ODE constraint acts as a scientific prior, with stronger dynamical enforcement improving external validation even as it reduces in-sample fit. Comparison with the standard EpiEstim method shows that our model-based approach substantially outperforms purely statistical $R_t$ estimation ($r=0.79$ vs.\ $r=0.25$) when applied to noisy case data.
\end{abstract}


%==========================================================================
\section{Introduction}\label{sec:intro}
%==========================================================================

Mechanistic models based on systems of ordinary differential equations (ODEs) are widely used to study infectious disease transmission and to inform public health decisions \citep{Anderson1991,Keeling2008,Hethcote2000}. A common starting point is the susceptible--exposed--infectious--removed (SEIR) model, which divides the population into four compartments---susceptible~($S$), exposed but not yet infectious~($E$), infectious~($I$), and removed~($R$)---and describes transitions between them over time. In its basic form, the SEIR model for population fractions reads
\begin{align}
  \dot{S}(t) &= -\beta(t)\,S(t)\,I(t), \label{eq:seir-S} \\
  \dot{E}(t) &= \beta(t)\,S(t)\,I(t) - \sigma E(t), \label{eq:seir-E} \\
  \dot{I}(t) &= \sigma E(t) - \gamma I(t), \label{eq:seir-I} \\
  \dot{R}(t) &= \gamma I(t), \label{eq:seir-R}
\end{align}
where $\sigma>0$ is the rate at which exposed individuals become infectious (mean latent period $1/\sigma$), $\gamma>0$ is the removal rate (mean infectious period $1/\gamma$), and $\beta(t)\geq 0$ is the time-varying transmission rate that captures changes in contact patterns, behavior, and policy interventions \citep{Keeling2008,Anderson1991,Flaxman2020}. The model conserves population mass: $S(t)+E(t)+I(t)+R(t)=1$ for all~$t$, and the basic reproduction number in a fully susceptible population is $R_0=\beta/\gamma$.

Calibrating such models to real surveillance data presents three interconnected challenges. First, the observed quantities---reported cases, hospital admissions, deaths---are not direct measurements of the SEIR compartments but rather indirect and noisy proxies, affected by reporting delays, incomplete ascertainment, changing case definitions, and evolving testing practices \citep{McGough2020,Gostic2020}. For instance, reported COVID-19 cases in the Netherlands increased not only because infections rose but also because testing capacity expanded from restricted hospital-based testing in March~2020 to community-wide testing by June~2020 \citep{Geubbels2023}. Second, when $\beta(t)$ is allowed to vary flexibly over time---as it must to capture the effects of lockdowns, behavioral changes, and seasonal forcing---the model becomes weakly identified: multiple combinations of $\beta(t)$ trajectories and scaling parameters can fit the observed data equally well \citep{Ramsay2007,QiZhao2010}. Third, when multiple data sources are available, they often disagree because of different reporting pipelines, case definitions, and post-processing steps \citep{McGough2020,Gostic2020}. A practical calibration method must be able to combine these heterogeneous streams while acknowledging and quantifying their relative information content.

Generalized profiling (GP), introduced by \citet{Ramsay2007}, offers an elegant approach to these challenges. The key idea is to represent the unknown state trajectories using a flexible basis expansion---typically B-splines---and to estimate the coefficients by minimizing an objective that balances fidelity to the data against compliance with the ODE system through a penalty term. This formulation has three important advantages over direct numerical integration of the ODEs. First, the B-spline trajectories are explicit polynomial functions of their coefficients, with derivatives obtained analytically by evaluating basis derivatives---no numerical ODE solver is embedded in the optimization loop, which greatly simplifies the computational structure. Second, the ODE enters as a soft penalty rather than a hard constraint, meaning the fitted trajectories can deviate from the deterministic SEIR dynamics when the data strongly disagrees; this is desirable because the SEIR model is only an approximation of the true epidemic process, which involves stochasticity, spatial heterogeneity, and behavioral feedbacks not captured by four deterministic compartments. Third, the approach naturally accommodates multiple data sources through separate observation operators, each linking a data stream to the appropriate SEIR compartment through an epidemiologically motivated mapping. GP has been shown to produce asymptotically efficient estimators under regularity conditions relating the spline resolution and penalty strength to the sample size \citep{QiZhao2010}, and has been applied to infectious disease dynamics by \citet{Hooker2011} for measles in Ontario. Practical implementations are available in the CollocInfer R package \citep{Hooker2016} and the pCODE package \citep{WangCao2023}.

In this paper, we develop a multi-source generalized profiling framework specifically designed for SEIR models with heterogeneous COVID-19 surveillance data. Our contributions are:
\begin{enumerate}
  \item We formulate nonlinear observation operators that link each of four data streams---reported cases, hospital admissions, ICU admissions, and digital contact monitoring data from the COVID~RADAR app---to the appropriate SEIR compartment, grounded in the epidemiology of the reporting process (\Cref{sec:obs-models}).
  \item We employ a log-state B-spline parameterization that guarantees positivity of all compartments without inequality constraints (\Cref{sec:logstate}), combined with a separate log-spline for the time-varying transmission rate (\Cref{sec:splines}), and parameter cascading that separates the approximately 120-dimensional trajectory estimation from the 21-dimensional transmission and scaling parameter inference (\Cref{sec:gp-objective}).
  \item We conduct a systematic source ablation study across all eight possible binary on/off configurations of the four data streams, quantifying the marginal information value of each source for recovering both the reproduction number and prevalence dynamics (\Cref{sec:results:ablation}).
  \item We perform a weight optimization over 64~configurations of per-source weights across all four streams, revealing the Pareto frontier between reproduction number and prevalence validation (\Cref{sec:results:weights}).
  \item We validate the model against two independent targets---RIVM reproduction number estimates and RIVM prevalence estimates---using comprehensive metrics including Pearson and Spearman correlations, RMSE, bias, and threshold agreement (\Cref{sec:results:validation}), and compare with the standard EpiEstim method of \citet{Cori2013} (\Cref{sec:results:epiestim}).
  \item We perform independent per-wave refitting to quantify how the constant detection fraction assumption breaks down across pandemic waves (\Cref{sec:results:waves}).
\end{enumerate}

The paper is organized as follows. \Cref{sec:background} reviews generalized profiling and related statistical approaches. \Cref{sec:method} presents the GP-SEIR formulation in detail. \Cref{sec:data} describes the Dutch COVID-19 data. \Cref{sec:results} reports numerical results across all analyses, \Cref{sec:future} discusses directions for future research, and \Cref{sec:conclusion} concludes.

%==========================================================================
\section{Background}\label{sec:background}
%==========================================================================

\subsection{Inference for ODE models from noisy and partial observations}

ODE models are used when a mechanistic description of system dynamics is available, but only noisy and indirect measurements of the system state are observed. In infectious disease surveillance, the latent epidemic process---the true numbers of susceptible, exposed, infectious, and removed individuals at each time point---is never directly observed. Instead, we observe quantities such as reported cases (a function of the exposed-to-infectious flow, filtered through testing capacity and reporting delays), hospital admissions (a fraction of infected individuals who become severely ill), and ICU admissions (a further subset). Each of these observed quantities is an indirect, noisy, and potentially biased proxy for the underlying SEIR state \citep{Gostic2020,McGough2020}. Moreover, different data products may disagree even when they aim to measure the same underlying quantity, because they are generated through different reporting pipelines with different definitions and processing steps.

Three broad classes of estimation approaches can be distinguished. The first class solves the ODE numerically (using, e.g., Runge--Kutta methods) and fits parameters by maximum likelihood or least squares. This approach is conceptually straightforward but computationally expensive because it requires a forward ODE solve at every iteration of the optimizer, plus sensitivity equations for gradient computation. The second class uses a two-stage method that first smooths the data nonparametrically and then matches the derivatives or moments of the smoothed curves to the ODE \citep{Liang2008}. This is computationally efficient but can lose statistical efficiency because the smoothing step does not account for the ODE structure. The third class---to which generalized profiling belongs---uses a hybrid approach that smooths the trajectories while simultaneously enforcing the ODE structure through a penalty term \citep{Ramsay2007,Hooker2016}. This retains the computational advantages of basis-function representations while incorporating the scientific knowledge encoded in the ODE.

\subsection{Generalized profiling (parameter cascading)}\label{sec:bg:gp}

Generalized profiling represents the unknown state trajectories using a basis expansion, typically B-splines:
\begin{equation}\label{eq:basis-expansion}
  x_j(t) \approx \sum_{k=1}^{K} c_{jk} B_k(t), \quad j \in \{S, E, I, R\},
\end{equation}
and estimates the coefficients $\{c_{jk}\}$ along with model parameters $\btheta$ by minimizing an objective that balances data fit and ODE compliance \citep{Ramsay2007}. The GP objective is
\begin{equation}\label{eq:gp-obj}
  \mathcal{J}(\bc,\btheta)
  = \underbrace{\sum_{m=1}^{M}\bigl\|\bm{y}_m - H_m(\bx(\cdot;\bc))\bigr\|^2_{W_m}}_{\text{data misfit}}
  + \underbrace{\eta\int\bigl\|\dot{\bx}(t;\bc)-f(\bx(t;\bc),\btheta,t)\bigr\|^2\dd t}_{\text{ODE penalty}}
  + \underbrace{\mathcal{R}(\bc,\btheta)}_{\text{regularization}},
\end{equation}
where $\eta>0$ controls the tradeoff between data fit and ODE adherence, $H_m$ are observation operators linking the $m$-th data stream to the state vector, $W_m$ are observation weights, and $\mathcal{R}$ collects additional regularization terms such as smoothness penalties.

The key computational advantage of GP is that the B-spline representation makes $\bx(t;\bc)$ and its derivative $\dot{\bx}(t;\bc)$ explicit functions of the coefficient vector $\bc$, evaluated by matrix multiplication with the basis matrices. No numerical ODE solver is needed: the ODE enters only through the penalty term, which compares the spline derivatives against the right-hand side of the ODE evaluated at the spline values.

\paragraph{Parameter cascading.}
The objective \eqref{eq:gp-obj} is typically minimized using a profiling (parameter cascading) strategy. The unknowns are split into two groups: the inner parameters $\bc$ (spline coefficients, high-dimensional) and the outer parameters $\btheta$ (model parameters, low-dimensional). For fixed $\btheta$, the inner problem for $\bc$ is a penalized nonlinear least-squares problem that can be solved efficiently using Levenberg--Marquardt or Gauss--Newton methods. Denoting the inner solution as $\hat{\bc}(\btheta)$, one defines the profiled objective $\widetilde{\mathcal{J}}(\btheta) := \mathcal{J}(\hat{\bc}(\btheta), \btheta)$ and optimizes over $\btheta$ in an outer loop. By the envelope theorem \citep{Ramsay2007}, the gradient of the profiled objective simplifies: at the inner optimum $\nabla_{\bc}\mathcal{J}=0$, so terms involving $\partial\hat{\bc}/\partial\btheta$ vanish, and $\nabla_{\btheta}\widetilde{\mathcal{J}} = \partial\mathcal{J}/\partial\btheta \big|_{\bc=\hat{\bc}(\btheta)}$.

\subsection{Statistical properties and related work}

\citet{QiZhao2010} analyzed the asymptotic efficiency and finite-sample properties of GP estimators, showing that they achieve the parametric rate under conditions that relate the spline approximation error, the penalty strength, and the sample size. This provides a theoretical foundation for understanding GP as more than an ad hoc regularizer: it is a statistically principled estimator when the smoothing and penalty parameters scale appropriately.

Several extensions and related methods have been developed. \citet{Brunel2015} connected GP to optimal control theory, providing an alternative derivation of the estimator through Pontryagin's maximum principle. \citet{Nanshan2022} extended the framework to non-Gaussian observations, which is relevant for count data in epidemiology. \citet{CaoRamsay2009} developed adaptive penalty selection via profiled generalized cross-validation, automating the choice of $\eta$. The pCODE R package \citep{WangCao2023} provides an alternative implementation with automatic Jacobian computation.

For epidemiological applications specifically, \citet{Hooker2011} demonstrated GP for measles transmission in Ontario, in a setting where the deterministic model is an imperfect representation of the true epidemic driven by stochastic dynamics and spatial heterogeneity. Their work showed that GP can recover meaningful transmission rate estimates even when the underlying ODE model is misspecified, a property that is crucial for our COVID-19 application.


%==========================================================================
\section{Method}\label{sec:method}
%==========================================================================

\subsection{SEIR dynamics and inferential target}\label{sec:seir}

We consider the SEIR model~\eqref{eq:seir-S}--\eqref{eq:seir-R} in terms of population fractions. Following standard practice in COVID-19 modeling, we treat the rate parameters $(\sigma, \gamma)$ as fixed from clinical literature and focus inference on the time-varying transmission rate $\beta(t)$ and the latent state trajectories $S(t), E(t), I(t), R(t)$ \citep{Keeling2008,Hethcote2000}. The rationale for fixing $(\sigma, \gamma)$ is that these parameters are well-characterized from clinical studies of viral shedding and symptom duration \citep{He2020,Byrne2020}, while $\beta(t)$ varies with behavioral and policy changes that cannot be determined a priori.

Allowing $\beta(t)$ to vary freely, however, creates an ill-posed inverse problem: for any smooth trajectory $I(t)$, one can always find a $\beta(t)$ that is consistent with the SEIR dynamics. In generalized profiling, this flexibility is controlled in two ways: (i)~the B-spline basis for $\beta(t)$ has limited resolution (knots every 28~days), preventing rapid oscillations; and (ii)~the ODE penalty ensures that the state trajectories are consistent with the SEIR dynamics at the chosen $\beta(t)$.

\subsection{Log-state parameterization}\label{sec:logstate}

SEIR compartments must satisfy $S, E, I, R > 0$ and $S+E+I+R=1$. Rather than imposing inequality constraints---which would complicate the optimization---we fit on the log scale. Define $z_j(t) = \log x_j(t)$, so that $x_j(t) = \exp(z_j(t)) > 0$ by construction, regardless of the value of $z_j$. Applying the chain rule $\dot{z}_j = \dot{x}_j / x_j$ to the SEIR system \eqref{eq:seir-S}--\eqref{eq:seir-R} yields the log-state ODE:
\begin{align}
  \dot{z}_S &= -\beta\,e^{z_I}, &
  \dot{z}_E &= \beta\,e^{z_S+z_I-z_E}-\sigma, \label{eq:log-SEIR}\\
  \dot{z}_I &= \sigma\,e^{z_E-z_I}-\gamma, &
  \dot{z}_R &= \gamma\,e^{z_I-z_R}. \nonumber
\end{align}
For example, the $\dot{z}_E$ equation follows from $\dot{z}_E = \dot{E}/E = (\beta SI - \sigma E)/E = \beta SI/E - \sigma = \beta \exp(z_S+z_I-z_E) - \sigma$. Similarly, $\beta(t) \geq 0$ is enforced by modeling $g(t) := \log\beta(t)$ and setting $\beta(t) = \exp\{g(t)\}$. This log-spline approach for transmission rates is standard in flexible SEIR analyses \citep{Frasso2016}.

\subsection{Spline representation}\label{sec:splines}

Each log-state is represented as a cubic B-spline (order~4) with knots placed every 14~days:
\begin{equation}
  z_j(t) \approx \sum_{k=1}^{K} c_{jk} B_k(t), \quad j \in \{S,E,I,R\},
\end{equation}
with derivatives obtained analytically: $\dot{z}_j(t) = \sum_k c_{jk}\dot{B}_k(t)$, where $\dot{B}_k$ is the derivative of the $k$-th basis function. The 14-day knot spacing is chosen to match the characteristic timescale of epidemic wave dynamics (roughly 2--4~weeks from onset to peak), while providing enough resolution to capture the rapid changes during lockdown transitions. Over the 366-day fitting window, this gives $K \approx 30$~basis functions per compartment and approximately $4K \approx 120$~inner unknowns.

The log-transmission uses a separate, coarser basis with knots every 28~days:
\begin{equation}
  g(t) = \log\beta(t) \approx \sum_{\ell=1}^{L} a_\ell b_\ell(t),
\end{equation}
where $L \approx 17$~basis functions. The coarser resolution reflects the fact that policy interventions and behavioral changes occur on slower timescales than the epidemic dynamics they produce. Together with four scaling parameters $(\rho, p_H, p_{\mathrm{ICU}}, \alpha_R)$, the outer problem has approximately $L+4 = 21$~unknowns.

\subsection{Observation models}\label{sec:obs-models}

Each data stream is linked to the latent SEIR states through an observation operator that reflects the epidemiology of how the data are generated. The choice of which SEIR compartment or flow to link each stream to is a modeling decision that should reflect the timing of the reporting process relative to the disease progression.

\paragraph{Cases $\to$ $E(t)$.}
Reported cases are modeled as proportional to the $E\to I$ flow:
\begin{equation}\label{eq:obs-cases}
  \mu_{\mathrm{case}}(t) = N\rho\sigma E(t),
\end{equation}
where $N=17{,}400{,}000$ is the Dutch population and $\rho\in(0,1]$ is a detection fraction. The choice to link cases to the $\sigma E$ flow (rather than to $\beta SI$ or $\gamma I$) is motivated by the timing of the diagnostic process: the flow $\sigma E$ represents symptom onset, which occurs approximately 5.5~days post-infection. Testing and reporting typically occur within 2--5~days of symptom onset, making the total delay from infection to reporting approximately 8--10~days \citep{Gostic2020}. Among the three SEIR flows---$\beta SI$ at infection (day~0), $\sigma E$ at symptom onset (day~5.5), and $\gamma I$ at removal (day~15)---the $\sigma E$ flow is temporally closest to the actual reporting event when explicit delay convolution is not used. The detection fraction $\rho$ absorbs the combined effects of under-ascertainment, testing capacity limitations, and the fraction of infections that are asymptomatic and never tested.

\paragraph{Hospital and ICU admissions $\to$ $I(t)$.}
Admissions are flows out of the infectious compartment, modeled as fixed fractions of the removal flow:
\begin{equation}\label{eq:obs-hosp}
  \mu_{\mathrm{hosp}}(t) = Np_H\gamma I(t), \qquad \mu_{\mathrm{ICU}}(t) = Np_{\mathrm{ICU}}\gamma I(t).
\end{equation}
This formulation assumes that hospital admission occurs at or near the time of removal from the infectious compartment. For COVID-19, the median time from symptom onset to hospitalization is approximately 5--7~days \citep{Dongelmans2022}, which is comparable to the remaining infectious period after symptom onset ($1/\gamma - 1/\sigma \approx 4$~days), making the removal flow a reasonable contemporaneous proxy. The parameters $p_H$ and $p_{\mathrm{ICU}}$ represent the probability that an infected individual is hospitalized or enters the ICU, respectively.

\paragraph{COVID~RADAR contact data $\to$ $I(t)/N$.}
The COVID~RADAR app, developed by the Leiden University Medical Center (LUMC) and the tech company ORTEC \citep{vanDijk2021,vanDijk2020}, was a questionnaire-based smartphone application active in the Netherlands from April~2020. Users self-reported daily information about their health, symptoms, and social contacts. Over $250{,}000$~individuals registered for the app, generating over 8.5~million observations \citep{SchootUiterkamp2025}, with approximately $10{,}000$~daily observations in steady state.

The key quantity we extract from the RADAR data is the ratio $c_I(t)/c(t)$, where $c(t)$ is the average daily number of contacts reported by users and $c_I(t)$ is the average number of those contacts who were later identified as infectious. Under the random mixing assumption of the SEIR model---which posits that each individual has an equal probability of contacting any other individual in the population---this ratio approximates the fraction of the population that is currently infectious \citep{SchootUiterkamp2025}:
\begin{equation}\label{eq:radar-theory}
  \frac{c_I(t)}{c(t)} \approx \frac{I(t)}{N}.
\end{equation}
Since $I(t)$ in our model is already expressed as a population fraction, the observation model simplifies to
\begin{equation}\label{eq:obs-radar}
  \mu_{\mathrm{radar}}(t) = \alpha_R I(t),
\end{equation}
where $\alpha_R < 1$ is a scaling parameter. The factor $\alpha_R$ corrects for the empirical observation, documented by \citet{SchootUiterkamp2025}, that RADAR-based prevalence estimates are systematically lower than RIVM prevalence estimates by a factor of approximately $0.41$ from October~2020 onward. This under-estimation likely reflects two biases: (i)~asymptomatic infectious individuals are less likely to be identified as infectious contacts because they do not report symptoms; and (ii)~the app's user base was not demographically representative of the general population, with under-representation of young ($<18$) and elderly ($>70$) age groups \citep{vanDijk2021}.

\paragraph{Fitting scale and weighting.}
Case, hospital, and ICU counts are fitted on the $\log(1+y)$ scale, which provides variance stabilization for count data spanning several orders of magnitude (daily cases range from 3 to nearly 13{,}000). RADAR fractions are fitted on the raw scale, as they are already approximately continuous and do not exhibit the extreme dynamic range of count data. Each data stream receives a base weight $W_m = 1/\mathrm{Var}(y_m^{\mathrm{transformed}})$, computed from the sample variance of the transformed observations, which normalizes streams to contribute roughly equally to the objective before per-source weight multipliers are applied.

\paragraph{Scaling parameters.}
The four scaling parameters $(\rho, p_H, p_{\mathrm{ICU}}, \alpha_R)$ must satisfy positivity and boundedness constraints. We enforce these through logistic transforms: each parameter $\theta_j$ is represented in the optimization as an unconstrained real number $u_j$, with $\theta_j = \theta_j^{\min} + (\theta_j^{\max} - \theta_j^{\min}) \cdot \text{logistic}(u_j)$. The bounds are: $\rho \in [0.05, 0.60]$, $p_H \in [0.002, 0.05]$, $p_{\mathrm{ICU}} \in [0.0003, 0.02]$, $\alpha_R \in [0.10, 1.00]$.

\subsection{Full objective and optimization}\label{sec:gp-objective}

Combining all components, the full GP-SEIR objective is:
\begin{equation}\label{eq:objective}
  \mathcal{J}(\bc,\ba) = \sum_{m=1}^{M} w_m \sum_{i=1}^{n_m}
    W_{m,i}\bigl(y_m(t_i)-H_m(\bz(t_i;\bc))\bigr)^2
  + \eta\int_{t_0}^{t_1}\bigl\|\dot{\bz}(t;\bc)-\tilde{f}(\bz(t;\bc),\ba,t)\bigr\|^2\dd t
  + \mathcal{R}(\bc,\ba),
\end{equation}
where $w_m \geq 0$ are per-source weights, $\tilde{f}$ is the log-state SEIR vector field \eqref{eq:log-SEIR}, and the integral is discretized on the daily time grid as $\sum_{q=1}^{Q} w_q \|\dot{\bz}(t_q;\bc) - \tilde{f}(\bz(t_q;\bc),\ba,t_q)\|^2$ with $w_q = 1$ (trapezoidal quadrature on daily data).

The regularization $\mathcal{R}$ includes three terms:
\begin{itemize}
  \item \emph{Mass conservation}: $\kappa_{\mathrm{mass}} \sum_q (S_q + E_q + I_q + R_q - 1)^2$ with $\kappa_{\mathrm{mass}} = 10^6$. This penalizes deviations from the constraint $S+E+I+R=1$, which the log-state parameterization does not automatically enforce.
  \item \emph{Roughness of $\log\beta(t)$}: $\kappa_\beta \int (g''(t))^2 \dd t$ with $\kappa_\beta = 10^3$. This prevents unrealistically rapid oscillations in the transmission rate.
  \item \emph{Soft upper bound on $\beta(t)$}: $\kappa_{\mathrm{cap}} \sum_q \max(\beta(t_q) - 0.6,\, 0)^2$ with $\kappa_{\mathrm{cap}} = 10^4$. This discourages biologically implausible transmission rates.
\end{itemize}

The ODE penalty weight $\eta = 10^6$ is the single most important tuning parameter. In log-state space, ODE residuals are typically of magnitude $10^{-2}$--$10^{-1}$, while data residuals on the $\log(1+y)$ scale are $0.5$--$2$. The weight $\eta = 10^6$ compensates for this scale mismatch, producing an effective ODE-to-data ratio of approximately $1{,}700\times$. The sensitivity analysis in \Cref{sec:results:sensitivity} demonstrates that this ratio is not arbitrary: it corresponds to the regime of optimal external validation.

The problem is non-convex due to the nonlinear observation models and exponential terms in the log-state ODE, so we use multi-start optimization with $N_{\mathrm{starts}}=5$ random initializations and retain the solution with the lowest objective value.

\paragraph{Parameter cascading implementation.}
The inner problem (approximately $4 \times 30 = 120$~state spline coefficients) is solved by Levenberg--Marquardt on the stacked residual vector. The outer problem (approximately 21~unknowns: 17~$\beta$-spline coefficients plus 4~scaling parameters) is solved by gradient descent with Armijo backtracking line search. The gradient is approximated by forward finite differences. Each outer iteration requires one inner solve for the objective and one inner solve per outer parameter for the gradient, giving approximately 22~inner solves per outer iteration. With 40~outer iterations, the total is approximately $40 \times 22 = 880$~inner solves per start, or $5 \times 880 = 4{,}400$~inner solves for 5~starts.

\subsection{Model-implied reproduction number}\label{sec:rt-method}

The effective reproduction number implied by the SEIR model is
\begin{equation}\label{eq:rt}
  R_t^{\mathrm{model}} = \frac{\beta(t)\,S(t)}{\gamma},
\end{equation}
which follows from the next-generation matrix approach \citep{Keeling2008}: in a fully susceptible population, $R_0 = \beta/\gamma$, and the effective number scales with the susceptible fraction. We compare this instantaneous reproduction number with RIVM's published $R_t$ estimates, which are computed using the Wallinga--Lipsitch method \citep{WallingaLipsitch2007}. The Wallinga--Lipsitch method relates the case reproduction number to the epidemic growth rate through the moment generating function of the generation interval distribution, and does not assume a specific compartmental structure. Importantly, it estimates a case reproduction number that can differ from the instantaneous $R_t$ in~\eqref{eq:rt}, particularly during periods of rapidly changing transmission \citep{Gostic2020,Vegvari2022}. Delay correction of RIVM's input data follows the P-spline nowcasting approach of \citet{vandeKassteele2019}. We set the weight of the RIVM $R_t$ series to zero in the fitting objective, ensuring the comparison constitutes genuine external validation.

%==========================================================================
\section{Case Study Data}\label{sec:data}
%==========================================================================

We apply the framework to the first year of the COVID-19 epidemic in the Netherlands (March~1, 2020 to March~1, 2021; $n=366$~days), a period spanning the first wave (March--June~2020), a summer inter-wave trough (June--October~2020), and the second wave (October~2020--March~2021). During this period, two national lockdowns were imposed (March~15--June~1, 2020 and October~14, 2020--end of study period), and testing capacity expanded substantially. \Cref{tab:data} summarizes all data streams.

\begin{table}[ht]
\centering\small
\caption{Data streams for the fitting window (March~1, 2020 to March~1, 2021; $n=366$~days). Fixed epidemiological parameters: $\sigma=1/5.5$~day$^{-1}$ (mean latent period 5.5~days), $\gamma=1/9.5$~day$^{-1}$ (mean infectious period 9.5~days) \citep{He2020,Byrne2020}, $N=17{,}400{,}000$.}
\label{tab:data}
\begin{tabular}{llcrll}
\toprule
Stream & Source & Days & Range & SEIR link & Observation model \\
\midrule
\multicolumn{6}{l}{\emph{Fitted streams}} \\
Cases & RIVM & 366 & 3--12{,}997/day & $\sigma E$ flow & $N\rho\sigma E$ \\
Hospital & NICE & 366 & 1--620/day & $\gamma I$ flow & $Np_H\gamma I$ \\
ICU & NICE & 366 & 0--127/day & $\gamma I$ flow & $Np_{\mathrm{ICU}}\gamma I$ \\
RADAR & DANS & 331 & 0.0002--0.005 & $I(t)/N$ & $\alpha_R I$ \\
\midrule
\multicolumn{6}{l}{\emph{Validation streams (not fitted)}} \\
$R_t$ & RIVM & 366 & 0.48--1.90 & $\beta S/\gamma$ & Wallinga--Lipsitch \\
Prevalence & RIVM & 366 & 3{,}334--168{,}222 & $\propto NI$ & Sero-surveys \\
\midrule
\multicolumn{6}{l}{\emph{Descriptive (not fitted, not validated)}} \\
Tests & RIVM & 274 & 94--80{,}581/day & --- & --- \\
Positivity & derived & 274 & 0.4\%--21.0\% & --- & --- \\
\bottomrule
\end{tabular}
\end{table}

\paragraph{Reported cases.} Daily reported cases are obtained from the RIVM municipality-level notification file, aggregated nationally \citep{Geubbels2023}. A total of $1{,}091{,}487$ cases were reported during the fitting window. Case counts are strongly influenced by testing capacity: during the first wave, testing was restricted to hospitalized patients, healthcare workers, and vulnerable groups, resulting in severe under-ascertainment; from June~2020, community-based testing became available to all symptomatic individuals, substantially increasing the detection rate.

\paragraph{Hospital and ICU admissions.} Daily hospital and ICU admissions are from the NICE registry (National Intensive Care Evaluation), accessed through the RIVM data portal. The raw ICU-to-hospital admission ratio across the full window is 18.4\%. Hospital admissions are considered more reliable than case counts because they are less affected by testing capacity changes, though they are subject to changes in admission criteria and the severity profile of infections across waves \citep{Dongelmans2022}. The ICU data is the noisiest stream, with 25~days of zero admissions.

\paragraph{COVID~RADAR contact data.} Following \citet{SchootUiterkamp2025}, we extract the 7-day smoothed ratio $c_I(t)/c(t)$ from the COVID~RADAR app data deposited at DANS \citep{vanDijk2020}. This ratio---the fraction of daily contacts that were with infectious individuals---serves as an estimate of $I(t)/N$ under SEIR random mixing (Equation~\eqref{eq:radar-theory}). The data are available for 331 of 366~days (starting April~5, 2020, due to the app launch date plus smoothing lag). Values range from $0.0002$ to $0.005$, meaning that between 0.02\% and 0.49\% of reported contacts were with infectious individuals at any given time.

\paragraph{RIVM reproduction number (validation only).} RIVM publishes daily $R_t$ estimates using the Wallinga--Lipsitch method \citep{WallingaLipsitch2007} with P-spline nowcasting \citep{vandeKassteele2019}. Importantly, the input data changed during our study window: before June~12, 2020, hospital admissions were used as the basis for $R_t$ estimation; afterward, GGD case notifications were used \citep{Geubbels2023}. This methodological change is relevant for interpreting the per-period validation results in \Cref{sec:results:validation}.

\paragraph{RIVM prevalence (validation only).} RIVM publishes daily estimates of the number of currently infectious individuals, derived from a combination of sero-surveys and hospitalization-based modeling \citep{Vos2021}. These estimates share no input data with any of the four fitted streams, making the prevalence correlation a fully independent validation metric.

\paragraph{Testing data (descriptive only).} Daily testing volumes are available from June~2020 onward (274~days). The positivity rate (positive tests / total tests) dropped from over 20\% in March~2020 to under 5\% by late summer, indicating a shift from highly targeted to broad-based testing. We use this data descriptively to contextualize changes in the effective detection rate.


%==========================================================================
\section{Results}\label{sec:results}
%==========================================================================

\subsection{Main fit}\label{sec:results:main}

Using parameter cascading with 5~multi-start initializations and equal source weights ($w_m=1$ for all four streams), the estimated scaling parameters are:
\begin{itemize}
  \item $\hat{\rho}=0.201$: approximately 20\% of infections are detected as reported cases. This is a time-averaged value; the per-wave analysis (\Cref{sec:results:waves}) shows it varies from $\sim$10\% to $\sim$35\%.
  \item $\hat{p}_H=0.012$: 1.2\% of infected individuals are hospitalized.
  \item $\hat{p}_{\mathrm{ICU}}=0.0025$: 0.25\% enter the ICU, giving $\hat{p}_{\mathrm{ICU}}/\hat{p}_H=20.8\%$.
  \item $\hat{\alpha}_R=0.406$: the RADAR contact fraction is approximately 41\% of the true prevalence fraction, closely matching the value of $\sim$0.41 reported by \citet{SchootUiterkamp2025} using an entirely independent methodology.
\end{itemize}

The fitted transmission rate $\hat{\beta}(t)$ ranges from $0.071$ to $0.233$~day$^{-1}$, and the model-implied reproduction number $R_t^{\mathrm{model}}=\hat{\beta}(t)\hat{S}(t)/\gamma$ ranges from $0.66$ to $2.20$. The susceptible fraction declines monotonically from $S(0)=0.995$ to $S(T)=0.844$, and the cumulative removed fraction reaches $R(T)=14.8\%$, corresponding to approximately $2.58$~million cumulative infections. This is consistent with Dutch seroprevalence surveys that found approximately 15\% of the population had been infected by early 2021 \citep{Vos2021}.

\begin{figure}[ht]
\centering
\begin{tikzpicture}
\begin{groupplot}[
  group style={group size=2 by 1, horizontal sep=1.5cm},
  width=0.48\textwidth, height=5cm, xlabel={Day}, grid=major,
  xmin=0, xmax=365, legend pos=north east, legend style={font=\scriptsize}
]
\nextgroupplot[ylabel={Daily cases}, title={A.\ Cases ($\hat\rho=0.201$)}]
  \addplot[only marks, mark=*, mark size=0.3, gray!60] table[x=day,y=cases_raw]{tikz_cases.csv};
  \addplot[colcases, thick, no markers] table[x=day,y=mu_cases]{tikz_cases.csv};
  \legend{Data, Model}
\nextgroupplot[ylabel={Admissions/day}, title={B.\ Hospital ($\hat p_H=0.012$)}]
  \addplot[only marks, mark=*, mark size=0.3, gray!60] table[x=day,y=hosp_raw]{tikz_hosp.csv};
  \addplot[colhosp, thick, no markers] table[x=day,y=mu_hosp]{tikz_hosp.csv};
\end{groupplot}
\end{tikzpicture}
\caption{Model fit to two of the four data streams. The model captures the overall dynamics of both waves but systematically over-predicts wave~1 cases (where constant $\hat\rho=0.201$ is too high) and under-predicts wave~2 cases (where it is too low).}\label{fig:main-fit}
\end{figure}

\subsection{Reproduction number validation}\label{sec:results:validation}

\Cref{tab:rt-metrics} provides a comprehensive summary of the $R_t$ validation. We report six complementary metrics: Pearson and Spearman correlations (measuring linear and rank agreement), $R^2$ (variance explained), RMSE and MAE (error magnitude), and mean bias (systematic over/under-estimation).

\begin{table}[ht]
\centering\small
\caption{Validation of model-implied $R_t^{\mathrm{model}}$ against RIVM estimates ($n=366$~days). RIVM used hospital admissions before June~12, 2020, and case notifications afterward.}
\label{tab:rt-metrics}
\begin{tabular}{lcccccc}
\toprule
Period & Days & Pearson & Spearman & RMSE & MAE & Bias \\
\midrule
Full period     & 366 & 0.824 & 0.826 & 0.161 & 0.128 & $+$0.027 \\
Pre-June 2020   & 103 & 0.894 & 0.563 & 0.162 & 0.125 & $-$0.055 \\
Post-June 2020  & 263 & 0.733 & 0.762 & 0.161 & 0.129 & $+$0.059 \\
\midrule
Wave~1 (Mar--Jun) &  92 & 0.926 & --- & 0.137 & --- & --- \\
Inter-wave (Jun--Oct) & 122 & 0.630 & --- & 0.213 & --- & --- \\
Wave~2 (Oct--Mar)  & 152 & 0.597 & --- & 0.122 & --- & --- \\
\bottomrule
\end{tabular}
\end{table}

The overall $R^2=0.536$ indicates the model explains approximately half of the variance in RIVM's $R_t$, with the remainder attributable to conceptual differences between instantaneous and case reproduction numbers \citep{Gostic2020,Vegvari2022}, reporting delays, and model approximation error. As a binary classifier for whether $R_t>1$ (epidemic growing) or $R_t<1$ (declining), the model agrees with RIVM on 75.4\% of days (136/167 growing days correctly identified, 140/199 declining days).

\begin{figure}[ht]
\centering
\begin{tikzpicture}
\begin{axis}[
  width=0.9\textwidth, height=5.5cm, xlabel={Day}, ylabel={$R_t$},
  xmin=0, xmax=365, ymin=0.3, ymax=2.4, grid=major, legend pos=north east, legend style={font=\scriptsize}
]
  \addplot[colrivm, thick, no markers] table[x=day,y=rt_rivm]{tikz_rt.csv};
  \addplot[colmodel, thick, no markers, dashed] table[x=day,y=rt_model]{tikz_rt.csv};
  \draw[gray, dashed] (axis cs:0,1) -- (axis cs:365,1);
  \legend{RIVM $R_t$ (not fitted), Model $R_t^{\mathrm{model}}$}
\end{axis}
\end{tikzpicture}
\caption{Reproduction number validation ($r=0.824$, $r_s=0.826$). The model captures the major dynamics---the first-wave peak, lockdown-induced decline, summer trough, and second-wave rise---but shows systematic bias that reverses sign between periods.}\label{fig:rt}
\end{figure}

The per-period decomposition reveals important patterns. The pre-June period shows high Pearson ($0.894$) but low Spearman ($0.563$): the strong monotonic decline from $R_t\approx 2.2$ to $R_t\approx 0.7$ during the first lockdown inflates Pearson (which captures the overall trend) while Spearman (which measures rank agreement of daily fluctuations) is low because the model's smooth B-spline trajectory cannot capture the rapid day-to-day variations in RIVM's estimates. After June, the pattern reverses: Pearson drops to $0.733$ while Spearman rises to $0.762$, indicating that the model captures the relative ordering of $R_t$ values better than their absolute levels. The sign of the bias also flips: $-0.055$ before June (model underestimates) and $+0.059$ after (model overestimates), suggesting a systematic shift related to the change in RIVM's $R_t$ input data.

\begin{figure}[ht]
\centering
\begin{tikzpicture}
\begin{axis}[
  width=0.9\textwidth, height=5.5cm, xlabel={Day}, ylabel={Infectious count}, scaled y ticks=false,
  xmin=0, xmax=365, grid=major, legend pos=north west, legend style={font=\scriptsize},
  yticklabel style={/pgf/number format/fixed}
]
  \addplot[colrivm, thick, no markers] table[x=day,y=prev_rivm]{tikz_prev.csv};
  \addplot[blue, thick, no markers, dashed] table[x=day,y=I_model]{tikz_prev.csv};
  \legend{RIVM prevalence (not fitted), Model $N\hat I(t)$}
\end{axis}
\end{tikzpicture}
\caption{Prevalence validation ($r=0.847$). RIVM prevalence estimates share no input data with any fitted stream, providing fully independent validation.}\label{fig:prev}
\end{figure}

\subsubsection{Validation independence}\label{sec:results:independence}

A concern with using RIVM's $R_t$ for validation is partial data circularity: after June~2020, RIVM computed $R_t$ from case notifications, which our model also fits. We address this with three lines of evidence. First, the per-period $R_t$ correlation is higher before June~12 ($r=0.894$) than after ($r=0.733$), despite the pre-June period using hospital-based RIVM input that shares data with our hospital stream. If circularity were driving the agreement, the case-based period should show higher correlation, not lower. Second, the ablation study (\Cref{sec:results:ablation}) shows that a model fitted entirely without case data (hosp+icu+radar) still achieves $R_t=0.807$ against the case-based RIVM $R_t$. Third, the prevalence validation ($r=0.847$) involves zero shared input data.

\subsection{Comparison with EpiEstim}\label{sec:results:epiestim}

To benchmark our model-based $R_t$ estimation, we compare with the widely used Cori method \citep{Cori2013,Thompson2019} as implemented in the EpiEstim R package. The Cori method models daily incidence $I_t$ as a Poisson renewal process: $I_t \sim \mathrm{Poisson}(R_t \sum_{s=1}^{t-1} I_{t-s}\, \omega_s)$, where $\omega_s$ is the discretized serial interval distribution. Under a conjugate Gamma prior, the posterior distribution of $R_t$ is available in closed form over user-specified sliding windows, making the method computationally fast and straightforward to implement. We use a 7-day window with a parametric serial interval (Gamma, mean 6.5~days, SD 3.5~days), consistent with COVID-19 estimates \citep{He2020}.

A fundamental distinction is that the Cori method was designed for symptom onset incidence. When applied to such data with known serial intervals, it yields exact $R_t$ estimates, albeit time-lagged by approximately the incubation period \citep{Cori2013}. In practice, however, symptom onset data is rarely available in real time. When applied to \emph{reported} case counts---as is common during outbreak response and as we do here---the method inherits all reporting artefacts: variable testing delays, day-of-week effects, changes in testing capacity, and administrative backlogs. These propagate directly into the $R_t$ estimates because the Poisson renewal model has no mechanism to distinguish true transmission changes from reporting noise.

Over the 358~days where all three $R_t$ estimates are available (\Cref{fig:epiestim}), our GP-SEIR model achieves $r=0.790$ against RIVM, while EpiEstim achieves only $r=0.245$. The EpiEstim estimates range from $0.56$ to $9.29$, with extreme values corresponding to artefactual case count spikes. In addition to the noise, there is a visible time lag in the EpiEstim estimates relative to RIVM. This lag has two sources: (i)~the convention of reporting $R_t$ at the end of the 7-day sliding window introduces a $\sim$3.5-day shift \citep{Gostic2020}; and (ii)~the use of reported cases rather than onset dates adds an additional $\sim$5--7 days of reporting delay. A lag analysis confirms this: the correlation between Cori and RIVM is maximized at a backward shift of approximately 10~days, reaching $r=0.68$. Even with this optimal shift, EpiEstim still substantially underperforms the GP-SEIR model ($r=0.79$ with no shift required), demonstrating that the ODE constraint provides fundamentally better recovery of transmission dynamics from noisy reported data.

\begin{figure}[ht]
\centering
\begin{tikzpicture}
\begin{axis}[
  width=0.9\textwidth, height=5.5cm, xlabel={Day}, ylabel={$R_t$},
  xmin=0, xmax=365, ymin=0, ymax=3.5, grid=major, legend pos=north east, legend style={font=\scriptsize}
]
  \addplot[colrivm, thick, no markers] table[x=day,y=rt_rivm]{tikz_epiestim.csv};
  \addplot[colmodel, thick, no markers, dashed] table[x=day,y=rt_model]{tikz_epiestim.csv};
  \addplot[colcori, thick, no markers, densely dotted] table[x=day,y=rt_cori_mean]{tikz_epiestim.csv};
  \draw[gray, dashed] (axis cs:0,1) -- (axis cs:365,1);
  \legend{RIVM, GP-SEIR ($r=0.79$), EpiEstim ($r=0.25$)}
\end{axis}
\end{tikzpicture}
\caption{Comparison of $R_t$ estimation methods. The GP-SEIR model substantially outperforms EpiEstim applied to raw reported cases. The visible lag in EpiEstim reflects both the 7-day window convention and unaccounted reporting delays. Even with an optimal 10-day backward shift ($r=0.68$), EpiEstim underperforms the unshifted GP-SEIR model. Values above $3.5$ are clipped.}\label{fig:epiestim}
\end{figure}

\subsection{Source ablation study}\label{sec:results:ablation}

We systematically include and exclude data streams to quantify their marginal information value, running all eight configurations with at least one stream. \Cref{tab:ablation} reports results sorted by prevalence correlation.

\begin{table}[ht]
\centering\small
\caption{Source ablation study. All configurations use the same penalties and 5~multi-start initializations.}
\label{tab:ablation}
\begin{tabular}{lccccccc}
\toprule
Sources & $R_t$ Pears. & $R_t$ RMSE & Prev Pears. & $\hat\rho$ & $\hat\alpha_R$ & $R(\text{end})$ \\
\midrule
hosp+icu+radar    & 0.807 & 0.176 & \textbf{0.895} & 0.201 & 0.410 & 15.9\% \\
cases+radar       & 0.822 & 0.161 & 0.880 & 0.201 & 0.406 & 15.5\% \\
hosp+icu          & 0.774 & 0.232 & 0.865 & 0.252 & 0.343 & 17.6\% \\
all four          & 0.710 & 0.230 & 0.862 & 0.204 & 0.416 & 16.1\% \\
radar only        & \textbf{0.823} & \textbf{0.156} & 0.826 & 0.201 & 0.408 & 20.3\% \\
cases+hosp+icu    & 0.695 & 0.260 & 0.812 & 0.179 & 0.395 & 21.0\% \\
cases+hosp        & 0.788 & 0.203 & 0.779 & 0.201 & 0.406 & 18.4\% \\
cases only        & 0.792 & 0.201 & 0.763 & 0.202 & 0.406 & 18.0\% \\
\bottomrule
\end{tabular}
\end{table}

Three central findings emerge. First, RADAR has the highest marginal value of any stream: adding RADAR to cases improves prevalence from $0.763$ to $0.880$ ($+0.117$), while adding hospital and ICU to cases yields only $+0.049$. Second, hosp+icu+radar without any case data achieves the best prevalence ($0.895$), better than the all-four configuration ($0.862$). Third, the ranking changes when sorted by $R_t$: radar-only achieves the highest $R_t$ correlation ($0.823$), while traditional surveillance without RADAR (cases+hosp+icu: $0.695$) performs worst. The $R_t$ and prevalence rankings are only weakly correlated ($r=0.10$), indicating these metrics capture complementary aspects of model quality.

\begin{figure}[ht]
\centering
\begin{tikzpicture}
\begin{axis}[
  width=0.9\textwidth, height=5.5cm, xlabel={Day}, ylabel={$\beta(t)$},
  xmin=0, xmax=365, ymin=0, ymax=0.35, grid=major, legend pos=north east, legend style={font=\tiny}
]
  \addplot[colcases, thick, no markers] table[x=day,y=cases_only]{tikz_beta_ablation.csv};
  \addplot[colradar, thick, no markers] table[x=day,y=radar_only]{tikz_beta_ablation.csv};
  \addplot[blue, thick, no markers, dashed] table[x=day,y=cases_radar]{tikz_beta_ablation.csv};
  \addplot[colhosp, thick, no markers, densely dotted] table[x=day,y=hosp_icu_radar]{tikz_beta_ablation.csv};
  \legend{Cases only, RADAR only, Cases+RADAR, Hosp+ICU+RADAR}
\end{axis}
\end{tikzpicture}
\caption{Transmission rate $\hat\beta(t)$ across source configurations. Configurations with RADAR produce smoother, more consistent $\beta(t)$ trajectories, while cases-only shows more volatility.}\label{fig:beta-ablation}
\end{figure}

\subsection{Weight optimization}\label{sec:results:weights}

We search over $4\times 2\times 2\times 4=64$~configurations: $w_{\mathrm{cases}}\in\{0.1,0.5,1,2\}$, $w_{\mathrm{hosp}}\in\{0,1\}$, $w_{\mathrm{ICU}}\in\{0,1\}$, $w_{\mathrm{radar}}\in\{0,0.5,1,2\}$. The best by prevalence: $(w_c{=}2,w_h{=}1,w_i{=}0,w_r{=}2)\to 0.908$; by $R_t$: $(0.1,0,0,0.5)\to 0.831$. A Pareto-efficient balance: $(1,1,0,2)\to R_t{=}0.813$, prev$=0.906$. All top-5 $R_t$ configurations have $w_{\mathrm{hosp}}=w_{\mathrm{ICU}}=0$. The correlation between $R_t$ and prevalence metrics across the 64~configurations is $r=0.32$---positive but modest---confirming they measure related but distinct aspects of model quality.

\subsection{Effective detection rate}\label{sec:results:detection}

We define $\rho_{\mathrm{eff}}(t) = y_{\mathrm{cases}}(t)/(N\sigma\hat{E}(t))$, the ratio of observed to model-predicted case counts. If constant $\rho$ were correct, $\rho_{\mathrm{eff}}(t)$ would be flat.

\begin{figure}[ht]
\centering
\begin{tikzpicture}
\begin{axis}[
  width=0.85\textwidth, height=5cm, xlabel={Day}, ylabel={$\rho_{\mathrm{eff}}(t)$},
  xmin=0, xmax=365, ymin=0, ymax=1, grid=major, legend pos=north west, legend style={font=\scriptsize}
]
  \addplot[blue, thick, no markers] table[x=day,y=rho_eff]{tikz_detection.csv};
  \addplot[red, thick, dashed, no markers] table[x=day,y=rho_const]{tikz_detection.csv};
  \legend{$\rho_{\mathrm{eff}}(t)$, Constant $\hat\rho=0.201$}
\end{axis}
\end{tikzpicture}
\caption{Effective detection rate. Wave~1 median: $0.15$; wave~2 median: $0.47$ ($3.2\times$ increase).}\label{fig:detection}
\end{figure}

The median $\rho_{\mathrm{eff}}$ is $0.15$ in wave~1 and $0.47$ in wave~2 ($3.2\times$), reflecting the expansion of testing capacity confirmed by the testing volume data (positivity rate dropped from $>$20\% to $<$5\%). This explains why reducing case weight improves external validation: with a single $\rho$, the model cannot simultaneously match wave~1 and wave~2 case counts.

\subsection{Penalty weight sensitivity}\label{sec:results:sensitivity}

\begin{table}[ht]
\centering\small
\caption{Penalty weight sensitivity ($\eta=10^6$ fixed, $W_{\mathrm{data}}$ varies).}
\label{tab:wdata}
\begin{tabular}{rrcccr}
\toprule
$W_{\mathrm{data}}$ & Ratio & $R_t$ & Prev & Cases RMSE & $R(\text{end})$ \\
\midrule
$10^{-2}$ & $17{,}343\times$ & \textbf{0.824} & 0.864 & 2816 & 15.6\% \\
$1$ & $1{,}734\times$ & 0.804 & 0.867 & 2836 & 15.2\% \\
$10^{2}$ & $173\times$ & 0.713 & \textbf{0.906} & 2430 & 22.5\% \\
$10^{4}$ & $17\times$ & 0.606 & 0.879 & 2074 & 29.2\% \\
$10^{6}$ & $1.7\times$ & 0.403 & 0.871 & 566 & 13.3\% \\
$10^{8}$ & $0.2\times$ & 0.432 & 0.871 & 537 & 6.2\% \\
\bottomrule
\end{tabular}
\end{table}

\begin{remark}[The ODE as a scientific prior]
$R_t$ correlation drops monotonically from $0.824$ to $0.403$ as data emphasis increases, while case RMSE drops from $2{,}816$ to $537$. At $W_{\mathrm{data}}=10^8$, $R(\mathrm{end})=6.2\%$---far below seroprevalence ($\sim$15\%)---showing that without ODE enforcement, the model loses epidemiological meaning. This confirms weak identifiability and demonstrates that the SEIR constraint acts as a scientific prior.
\end{remark}

\subsection{Residual diagnostics}\label{sec:results:residuals}

\begin{figure}[ht]
\centering
\begin{tikzpicture}
\begin{groupplot}[
  group style={group size=2 by 1, horizontal sep=1.5cm},
  width=0.48\textwidth, height=4.5cm, grid=major
]
\nextgroupplot[xlabel={Lag (days)}, ylabel={ACF}, title={Case residual ACF}, ymin=-0.1, ymax=1]
  \addplot[blue, thick, mark=*, mark size=1.5] table[x=lag,y=acf]{tikz_acf.csv};
  \draw[red, dashed] (axis cs:0,0.104) -- (axis cs:30,0.104);
  \draw[red, dashed] (axis cs:0,-0.104) -- (axis cs:30,-0.104);
\nextgroupplot[xlabel={Theoretical quantile}, ylabel={Sample quantile}, title={Q-Q plot}]
  \addplot[only marks, mark=*, mark size=0.5, gray!60] table[x=theoretical,y=sample]{tikz_qq.csv};
  \addplot[red, thick, no markers, domain=-3:3]{x};
\end{groupplot}
\end{tikzpicture}
\caption{Case residuals on $\log(1+y)$ scale show strong autocorrelation (ACF lag-1: $0.94$) and a heavy left tail (Q-Q min: $-5.6$), both consistent with the constant-$\rho$ misspecification.}\label{fig:residuals}
\end{figure}

\subsection{Per-wave refitting}\label{sec:results:waves}

\begin{table}[ht]
\centering\small
\caption{Per-wave results: independent refit (separate B-spline bases) vs.\ main fit sliced to each period.}
\label{tab:wave-refit}
\begin{tabular}{lcccccc}
\toprule
Period & $\hat\rho$ & $\hat p_H$ & $\hat\alpha_R$ & $R_t$ corr & Prev corr & Det.\ rate \\
\midrule
\multicolumn{7}{l}{\emph{Independent refit}} \\
Wave~1 (Mar--Jun) & 0.097 & 0.016 & 0.614 & 0.873 & 0.630 & 1.1\% \\
Inter-wave (Jun--Oct) & 0.251 & 0.010 & 0.567 & 0.493 & 0.976 & 2.6\% \\
Wave~2 (Oct--Mar) & 0.350 & 0.012 & 0.286 & 0.635 & 0.617 & 5.1\% \\
\midrule
\multicolumn{7}{l}{\emph{Sliced from main fit ($\hat\rho=0.201$)}} \\
Wave~1 & --- & --- & --- & 0.925 & 0.907 & --- \\
Inter-wave & --- & --- & --- & 0.630 & 0.973 & --- \\
Wave~2 & --- & --- & --- & 0.597 & 0.263 & --- \\
\bottomrule
\end{tabular}
\end{table}

The detection fraction increases 3.6-fold from $\hat\rho=0.097$ (wave~1) to $0.350$ (wave~2), directly quantifying the impact of testing expansion. The sliced wave-2 prevalence ($0.263$) is strikingly low because constant $\hat\rho=0.201$ underestimates wave-2 detection, distorting the $E(t)$ trajectory.

\subsection{Multi-start uncertainty}\label{sec:results:multistart}

Across 5~starts: $\hat\rho\in[0.131,0.253]$, $\hat\alpha_R\in[0.381,0.439]$, $R_t$ corr $\in[0.689,0.827]$, prev corr $\in[0.800,0.853]$. The best two starts give nearly identical validation ($R_t$: $0.824$ vs.\ $0.825$) despite different $\hat\rho$ ($0.201$ vs.\ $0.210$), suggesting a ridge in the objective landscape.


%==========================================================================
\section{Future Research Directions}\label{sec:future}
%==========================================================================

Several promising extensions emerge from this work.

\paragraph{Time-varying detection fraction $\rho(t)$.} The per-wave analysis shows $\hat\rho$ varies 3.6$\times$ across the study period. Modeling $\rho(t)$ as a separate B-spline with its own roughness penalty would eliminate the dominant source of misspecification (case residual ACF lag-1: $0.94$) and substantially improve prevalence validation. Preliminary experiments with $\rho(t)$ suggest that careful tuning of the roughness penalty is needed to prevent the optimizer from absorbing all model-data discrepancy into $\rho(t)$ rather than adjusting $\beta(t)$. Incorporating testing volume data as an informative prior for $\rho(t)$ could provide the necessary regularization.

\paragraph{Delay convolution operators.} Cases are reported $\sim$8--10~days after infection, but our model uses $\sigma E$ (day~5.5) as a contemporaneous proxy. Replacing $N\rho\sigma E(t)$ with $N\rho \sum_\tau h(\tau)\sigma E(t-\tau)$ for a fixed delay kernel $h$ (e.g., log-normal with mean 3~days, capturing the test-to-report delay) would improve temporal alignment. Similarly, hospital admissions could be convolved with a symptom-to-admission delay distribution.

\paragraph{Bayesian uncertainty quantification.} Our multi-start analysis provides a crude uncertainty range but not proper confidence intervals. Block bootstrap---resampling 7-day blocks of residuals to preserve the strong autocorrelation structure, adding them back to the fitted curves, and refitting---would provide pointwise confidence bands on $\beta(t)$, $R_t$, and all parameters. Alternatively, a Bayesian formulation with Hamiltonian Monte Carlo sampling of the outer parameters (treating the inner problem as a likelihood evaluation) could provide posterior distributions.

\paragraph{Forecasting.} The GP-SEIR framework naturally supports forecasting by extrapolating the B-spline bases beyond the fitting window. Combining the fitted state estimates at the end of the training period with scenario-based $\beta(t)$ projections (e.g., constant, linearly declining, or seasonal) could provide short-term forecasts of cases, hospitalizations, and prevalence.

\paragraph{Additional data sources.} Wastewater surveillance data, which became available in the Netherlands from September~2020 \citep{Geubbels2023}, could provide another independent prevalence proxy through an observation model linking viral RNA concentration in sewage to $I(t)$ via a shedding profile convolution. Contact survey data from the CoMix study \citep{Geubbels2023} could inform the mixing assumptions underlying the RADAR observation model.

\paragraph{Spatial and age-structured extensions.} The current model treats the Netherlands as a single well-mixed population. Extending to a metapopulation structure (e.g., by province) or an age-structured model would allow incorporating age-specific contact patterns and the documented shift toward younger demographics in the second wave \citep{Dongelmans2022}.


%==========================================================================
\section{Conclusion}\label{sec:conclusion}
%==========================================================================

We have developed and applied a multi-source generalized profiling framework for SEIR models with time-varying transmission. The method uses log-state B-splines for guaranteed positivity, parameter cascading for structured optimization, and epidemiologically motivated observation operators for linking heterogeneous surveillance data to latent epidemic compartments.

Applied to four Dutch COVID-19 data streams over the first pandemic year, the method produces several notable findings. The estimated RADAR scaling parameter $\hat\alpha_R=0.406$ independently matches the $\sim$0.41 reported by \citet{SchootUiterkamp2025}. The reproduction number correlates at $r=0.82$ (Spearman $r_s=0.83$, $R^2=0.54$, threshold agreement 75\%) with RIVM estimates, substantially outperforming EpiEstim ($r=0.25$). The prevalence validation achieves $r=0.85$ with zero shared input data. The source ablation reveals RADAR provides the highest marginal value, and hosp+icu+radar without cases achieves the best prevalence ($0.895$)---better than all four sources together ($0.862$). The weight optimization over 64~configurations confirms RADAR benefits from higher weight and ICU is optimally excluded. The per-wave analysis quantifies how constant $\hat\rho$ varies 3.6$\times$ across waves, explaining why lower case weight improves results and motivating time-varying $\rho(t)$ as the most promising extension.

\bigskip
\noindent\textbf{Acknowledgments.}
COVID~RADAR data: DANS (\doi{10.17026/dans-zcd-m9dh}).
RIVM data: \url{https://data.rivm.nl/covid-19/}.
Code: \url{https://github.com/hadiabp/gp-seir-multisource}.

\bibliographystyle{plainnat}
\bibliography{references}

\end{document}
